\section{Transient output faces and leakage-sensitive memory}\label{sec:r70-leakage}
Uniform decay of a physical state and decay of its observable probabilities do not measure the same obstruction to exact stochastic realization. In this section leakage removes population from a protected subspace, whereas internal depolarization changes the spectrum relative to that surviving population. We prove that the first operation costs exactly one additional live label in the nondegenerate single-frequency family, independently of how rapid the leakage is. We then obtain finite-model witnesses on a fixed six-dimensional squared-prediction experiment, and an exact finite-realization result when a positive background is added to the terminal effect.

Nonnegative realizations, dominant-pole obstructions and transient-state reductions have classical antecedents \cite{BenvenutiFarina04,OCinneide90,Maier91,MonrasWinter14}. We do not claim to originate these methods. The statements below specify the constrained stochastic decoder, prescribed initial label, sharp label increment, physical leakage channel and finite squared-task tests. All necessary positive-matrix arguments are included.

\subsection{The output-visible transient core}
For a nonzero sequence $u=(u_n)_{n\ge0}$ in $[0,1]$ tending to zero, let $s(u)$ be its least stochastic realization size:
\[
 u_n=e_1^{\mathsf T}Q^nd,\qquad Q\ge0,\quad Q\one=\one,\quad 0\le d\le\one.
\]
The initial label is fixed, not a freely randomized initial distribution. Define $t(u)$ by the same formula with a nonnegative row-substochastic matrix $S$, $S\one\le\one$, whose spectral radius is strictly below one. Both minima are infinite if no finite realization exists.

\begin{lemma}[Exact transient-core reduction]\label{lem:r70-core}
For every such sequence,
\begin{equation}\label{eq:r70-core-count}
 s(u)=t(u)+1.
\end{equation}
More precisely, any $M$-label stochastic realization has a nonempty set of at most $M-1$ reachable labels that can eventually give positive output. Its restriction $S$ realizes $u$, is transient and uses the same initial label. If also $u_n\le C\lambda^n$ for $0<\lambda<1$, then the spectral radius of this restriction is at most $\lambda$.
\end{lemma}
\begin{proof}
Remove labels not reachable from label one. Call a remaining label active if a positive-probability path from it reaches some label $j$ with $d_j>0$. The initial label is active because $u$ is nonzero: there is some $n$ with $e_1^{\mathsf T}Q^nd>0$, so a positive-weight path from label one reaches a label with positive decoder output. This remains true when $u_0=0$. A path that leaves the active set cannot return: a return would make every intervening reachable label active. Thus, with $S$ the active principal submatrix and $d_A$ its decoder, $u_n=e_1^{\mathsf T}S^nd_A$.

For each active $i$, choose a path of length $a_i$ from the start to $i$ of weight $p_i>0$. For each active $j$, choose a path of length $b_j$ from $j$ with positive terminal expected output $r_j>0$. Positivity gives, for every $n$,
\begin{equation}\label{eq:r70-path-bound}
 p_ir_j(S^n)_{ij}\le u_{a_i+n+b_j}.
\end{equation}
Consequently every entry of $S^n$ tends to zero. In finite dimension this forces every eigenvalue to have modulus less than one. In the exponential case it gives $(S^n)_{ij}\le C_{ij}\lambda^n$, and applying this bound to an eigenvector proves the stronger spectral-radius bound. There must be a reachable inactive label: otherwise $S$ would be stochastic and have eigenvalue one. Hence the active size is at most $M-1$, proving $t(u)+1\le s(u)$.

Conversely a $k$-label transient realization extends to the stochastic matrix and decoder
\begin{equation}\label{eq:r70-cemetery}
 \widehat Q=\begin{pmatrix}S&\one-S\one\\0&1\end{pmatrix},
 \qquad \widehat d=\binom{d_A}{0}.
\end{equation}
They realize the same sequence from the same initial label with $k+1$ labels. This proves the reverse inequality, including the infinite case. The added label is a live zero-output label until the actual herald; it is not permission to stop at an unobserved physical leakage event.
\end{proof}

We recall a small positive-matrix fact with its proof. If $B\ge0$ is irreducible and has a nonzero eigenvalue, it has a vector $h>0$ with $Bh=Rh$, where $R$ is its spectral radius. Indeed, for $B+\epsilon\one\one^{\mathsf T}$ the map sending a simplex vector to its normalized image has a fixed point by the finite-dimensional simplex fixed-point theorem. The fixed point is strictly positive. Along $\epsilon\downarrow0$ a convergent subsequence gives $Bh=Rh$, $h\ge0$, $\sum_i h_i=1$. If a coordinate vanished, the eigenvector equation and irreducibility would propagate that zero to every coordinate. Thus $h>0$, and $R>0$. The similarity $R^{-1}\operatorname{diag}(h)^{-1}B\operatorname{diag}(h)$ is stochastic, so its maximum-row-sum bound shows that $R$ is indeed the spectral radius. This is the Perron normalization used below, not a new Perron theorem.

\begin{proposition}[A relative spectral obstruction for decaying outputs]\label{prop:r70-relative-spectrum}
Suppose an $M$-label stochastic response satisfies $0\le u_n\le C\lambda^n$ with $0<\lambda<1$. If its rational generating function has a nonreal pole at $\zeta^{-1}$, then
\begin{equation}\label{eq:r70-relative-spectrum}
 \zeta/\lambda\in\mathcal K_{M-1}.
\end{equation}
Here $\mathcal K_j$ is the stochastic eigenvalue region defined in Section~\ref{sec:r69-spectral}. It is the decay-normalized mode, not merely the absolute mode $\zeta$, that constrains the active state count.
\end{proposition}
\begin{proof}
The core in Lemma~\ref{lem:r70-core} gives
$\sum_nu_nz^n=e_1^{\mathsf T}(I-zS)^{-1}d_A$. Hence $\zeta$ is an eigenvalue of $S$. Order its strongly connected components to obtain a block triangular matrix. One irreducible diagonal block $B$ has that eigenvalue. Its spectral radius $R$ obeys $|\zeta|\le R\le\lambda$. Perron normalization gives a stochastic matrix $T$ of size $l\le M-1$ with eigenvalue $\zeta/R$ and eigenvector $v$.

A stationary probability row $\pi$ for $T$ exists: take a convergent subsequence of the averages of any probability row and its first $n-1$ iterates. Since $\zeta/R\ne1$, one has $\pi v=0$. Put $a=R/\lambda\in(0,1]$. The stochastic matrix
$T'=aT+(1-a)\one\pi$ satisfies $T'v=(\zeta/\lambda)v$. Padding by unreachable self-loop labels, if needed, proves membership in $\mathcal K_{M-1}$. No global Perron scaling of the reducible core has been assumed.
\end{proof}

The last qualification matters. A reducible matrix such as
$S=\left(\begin{smallmatrix}1/2&1/4\\0&1/4\end{smallmatrix}\right)$
has spectral radius $1/2$, but no strictly positive $h$ satisfies $Sh\le h/2$. The proof normalizes the irreducible block carrying the observed pole, rather than silently imposing that false global inequality.

\subsection{An exact leakage law for every damping and phase}
Fix $0<\lambda<1$, $0<\alpha\le1$, $0<\gamma\le1$, and $\sin\vartheta\ne0$. Let
\begin{equation}\label{eq:r70-leaky-response}
 u_n=\frac{\lambda^n}{2}
       [1+\gamma\alpha^n\cos(n\vartheta+\varphi)].
\end{equation}
The scalar $\lambda$ is survival against leakage; $\alpha$ is coherence retention within the surviving qubit. These parameters are known properties of the prepared target channel, not extra observations available to the observer.

\begin{theorem}[Exact leakage increment]\label{thm:r70-leakage}
For the nondegenerate family \eqref{eq:r70-leaky-response},
\begin{equation}\label{eq:r70-exact}
 s(u)=1+m(\alpha,\vartheta)
     =1+\min\{k:\alpha e^{i\vartheta}\in\mathcal K_k\}.
\end{equation}
This holds over all stochastic matrices and randomized decoder rows, with a prescribed initial label. In particular the value is independent of every $0<\lambda<1$, every nonzero contrast and every preparation phase. For $\alpha<1$ it is finite. At $\alpha=1$ it is $h+1$ if $e^{i\vartheta}$ has finite order $h$, and it is infinite otherwise.

The full two-output Hankel response has rank four, regardless of this possibly much larger or infinite exact memory. In the fixed squared-task embedding it is the least live memory at the original cost $1/p$, with event-stopped and state-stopped total-label upper counts $\max\{s(u),2\}$ and $s(u)+2$ when finite.
\end{theorem}
\begin{proof}
The three modes $\lambda$, $\lambda\alpha e^{i\vartheta}$ and $\lambda\alpha e^{-i\vartheta}$ in $u$ are distinct and have nonzero coefficients because $\lambda>0$, $\alpha>0$, $\gamma>0$ and $\sin\vartheta\ne0$. The phase factor $e^{i\varphi}$ never vanishes, including at initial-output zeros. Thus its generating function has a pole at $(\lambda\alpha e^{i\vartheta})^{-1}$. Proposition~\ref{prop:r70-relative-spectrum} makes every finite realizing size $M$ satisfy $\alpha e^{i\vartheta}\in\mathcal K_{M-1}$. Theorem~\ref{thm:r69-spectral} proves the required lower and the infinite case.

For the upper take a minimal realization $(Q,d)$ of the unattenuated response
$b_n=[1+\gamma\alpha^n\cos(n\vartheta+\varphi)]/2$.
Replace its continuing block by $\lambda Q$ and append a live zero-output label, as in \eqref{eq:r70-cemetery}. Its output is exactly $\lambda^nb_n$ from time zero. This proves equality without an extra initialization operation. The unit-circle alternatives follow from the already proved root-of-unity characterization.

The scalar response spans its three exponential modes. Its complement adds the independent constant mode one, because $\lambda<1$. Vandermonde independence therefore gives full response rank four. Theorem~\ref{thm:r67-embedding} forces this entire probability response in the actual squared task; the role-packing statement is then inherited without identifying protected physical states with live observer labels.
\end{proof}

\begin{corollary}[Noise laws depend on relative damping]\label{cor:r70-rates}
For fixed $0<\lambda<1$ and $0<\alpha<1$, the worst-phase exact count is
\[
 \max_{\sin\vartheta\ne0}s(u)
 =1+\max\{3,\lceil\pi/\arccos\alpha\rceil\}.
\]
For every fixed phase with the global badly-approximable certificate of Theorem~\ref{thm:r69-cubic}, $s(u)-1$ has the same matching order $(1-\alpha)^{-1/3}$. Neither rate depends on leakage $\lambda$, even when leakage is allowed to vary simultaneously with $\alpha$. At the original algebraic phase, the first nontrivial threshold is exactly four live labels for $0<\alpha\le5/(3+4\sqrt3)$, and at least five beyond it. No badly-approximable claim is made for that algebraic phase.
\end{corollary}
\begin{proof}
Substitute \eqref{eq:r70-exact} into the exact worst-phase, fixed-phase and three-label results. All subtractions concern live labels, including the necessary zero-output label; no terminal label has been subtracted from an event-stopped dictionary.
\end{proof}

\subsection{A strictly contractive quantum channel with infinite exact memory}
Let $K=\mathbb C^2\oplus\mathbb C|\bot\rangle$, let $P$ project onto the qubit and let $\Omega=|\bot\rangle\langle\bot|$. With the qubit channel $\Phi_\alpha$ and preparation from Section~\ref{sec:r69-spectral}, define
\begin{equation}\label{eq:r70-channel}
 \Psi_{\lambda,\alpha}(X)
 =\lambda\Phi_\alpha(PXP)
  +[\tr((I-P)X)+(1-\lambda)\tr(PX)]\Omega.
\end{equation}
The qubit output is embedded in its original block and cross-block coherences are discarded. Prepare the qubit and audit the qubit effect $E=(I+\sigma_z)/2$, extended by zero at $\bot$.

\begin{theorem}[Uniform physical mixing versus exact response memory]\label{thm:r70-physical}
Equation \eqref{eq:r70-channel} is a completely positive trace-preserving channel with unique stationary state $\Omega$. For all density matrices $\rho,\sigma$ and all $n\ge0$,
\[
 \|\Psi_{\lambda,\alpha}(\rho)-\Psi_{\lambda,\alpha}(\sigma)\|_1
 \le\lambda\|\rho-\sigma\|_1,
 \qquad
 \|\Psi_{\lambda,\alpha}^n(\rho)-\Omega\|_1\le2\lambda^n.
\]
Its prepared audit response is exactly \eqref{eq:r70-leaky-response}, so its exact stochastic observer memory is \eqref{eq:r70-exact}. In particular, at $\alpha=1$ and the algebraic rotation $(3+4i)/5$, every $0<\lambda<1$ gives infinite exact memory despite strict trace-norm contraction and exponential convergence to one stationary state. For algebraic $\lambda$ and algebraic preparation data, all primitive superoperator entries are algebraic.

The squared-prediction embedding has protected dimension six, fixed dictionary $\{-1,1\}$, three actual audit outcomes and least expected cost $1/p$. In the specified irrational, internally undamped case, every finite live-memory class has an attained strictly larger risk than the full-history value $2$. A prefix with $T$ exact-output labels and one live zero-output tail label has the same cost and excess risk at most
\begin{equation}\label{eq:r70-approx}
 \frac{2\varepsilon p(q\bar t\lambda)^T}{1-q\bar t\lambda},
 \qquad q=1-p,\quad\varepsilon=1/4.
\end{equation}
\end{theorem}
\begin{proof}
Write $\Xi(X)=\Phi_\alpha(PXP)+\tr((I-P)X)\Omega$ and
$\mathcal E(X)=\tr(X)\Omega$. Both are channels: compression, the declared qubit channel and scalar replacement are completely positive, and their output traces sum to the input trace. Then $\Psi_{\lambda,\alpha}=\lambda\Xi+(1-\lambda)\mathcal E$. On a traceless Hermitian difference the second term vanishes. Trace-norm contraction of a channel follows directly by applying positivity and trace preservation to the positive and negative parts of that difference. This proves the contraction inequality. Since $\Omega$ is fixed, iteration gives the displayed bound and uniqueness.

Starting in the qubit, its surviving block after $n$ calls is $\lambda^n\Phi_\alpha^n(\rho)$; the complementary mass is at $\bot$. Taking the effect gives \eqref{eq:r70-leaky-response}. The preceding theorem proves the memory claim. At $(3+4i)/5$ the unit phase is not a root of unity, by the rational algebraic-integer argument already proved. The spectral obstruction is in the decaying output-visible block, not at the peripheral spectrum of the entire trace-preserving channel.

The embedding of Theorem~\ref{thm:r67-embedding} doubles the protected dimension and retains its real squared audit. Its unique optimal probability kernel must reproduce $u$. Compactness of each fixed-memory controller class and the uniform geometric word tail give attainment; equality at a finite size would contradict the exact-memory theorem. For the approximation, store $u_0,\ldots,u_{T-1}$ and thereafter output probability zero. Since $u_n\le\lambda^n$, summing the exact risk identity \eqref{eq:r67-generating-risk} yields \eqref{eq:r70-approx}. Nothing in this construction observes leakage or changes first-herald stopping.
\end{proof}

This is not the earlier unique-sign counterexample under depolarization. Here the required response is probability valued and its absolute nonconstant eigenvalues all lie strictly inside the unit disk. The impossibility comes from their phases relative to the visible decay rate. The stationary state is on a face, so the interior hypothesis of the instrument-wide realization theorem is genuinely unavailable. With $\gamma=1$ and $\varphi=\pi$, moreover, $u_0=0$ and $0\le u_n\le\lambda$ for all $n\ge1$. The one-label zero response is uniformly within $\lambda$ of this law while exact memory remains infinite for every positive $\lambda$. Exact realization is not a robust approximate-memory assertion.

\subsection{Finitely many physical models detect the obstruction}
The continuum of herald-flag parameters used to force an arbitrary full response can be replaced, against a prescribed finite-memory bound, by a finite set whose cardinality depends only on that bound and the target's linear order.

\begin{theorem}[Finite rational-resolvent tests and fixed-dimensional witnesses]\label{thm:r70-finite-tests}
Suppose a target channel response has generating function $U(z)=P(z)/D(z)$ with
$\deg P\le r-1$, $\deg D\le r$, and $D(z)\ne0$ on $[0,q\bar t]$. Fix $M\ge1$ and any $M+r$ distinct parameters $t_j\in[0,\bar t]$. In the squared embedding, an observer with at most $M$ live labels has risk exactly $2$ on the finite family $(t_j,\sigma)$, $\sigma\in\{-1,1\}$, if and only if it realizes the entire target response at every failure count.

Consequently for the algebraic leaking qubit with $\lambda=1/2$, $\alpha=\gamma=1$, $e^{i\vartheta}=(3+4i)/5$ and $\varphi=0$, the following holds for every $M$: a six-dimensional prepared squared experiment with only $2(M+3)$ rationally specified models, one unit-cost command and the fixed two-value dictionary has boundary value $2$, but every at-most-$M$-label observer has risk at least $2+\delta_M$, where $\delta_M>0$ is an attained, computable real algebraic number. One may take $p=1/2$, $\bar t=1/2$ and $t_j=j/[2(M+2)]$, $0\le j\le M+2$. Neither the physical dimension nor the target response order increases with $M$.
\end{theorem}
\begin{proof}
Boundary feasibility still forces first-herald commitment in this finite model family, since its physical herald probability is always $p$. An $M$-label rule has
$B(z)=e_1^{\mathsf T}(I-zQ)^{-1}d=P_Q(z)/D_Q(z)$,
where $\deg P_Q\le M-1$, $\deg D_Q\le M$ and $D_Q(z)\ne0$ for $0\le z<1$. Smaller matrices may be padded by unreachable rows.

The two signs and \eqref{eq:r67-generating-risk} make risk $2$ equivalent to $B(qt_j)=U(qt_j)$ at every chosen parameter. The polynomial
$P_QD-PD_Q$ has degree at most $M+r-1$ and vanishes at $M+r$ distinct points. It is therefore zero. Equality of the power series about zero gives every response coefficient. The converse follows immediately from the same risk identity. This is rational interpolation with known degree, not extrapolation from arbitrary fitted samples.

For the stated target,
\[
 U(z)=\frac1{2(1-\lambda z)}+
       \frac{1-\lambda cz}{2(1-2\lambda cz+\lambda^2z^2)},
 \qquad c=3/5,
\]
so $r=3$ suffices and all coefficients are rational. The exact leakage theorem excludes every finite realization. Hence the finite tests exclude equality for the specified $M$. The full-history target kernel still attains $2$ in every selected model. Compact controller simplexes and the geometric tail imply the fixed-$M$ minimum is attained, so its gap $\delta_M$ is strictly positive.

Finally introduce real variables $Q,d$ and vectors $v_j$ satisfying
$(I-qt_jQ)v_j=d$. These matrices are invertible for every stochastic $Q$, and the achieved risk is
$2+2\varepsilon p\max_j|(v_j)_1-U(qt_j)|$.
All defining constraints are polynomial over the rationals. Real-algebraic elimination computes the attained minimum and an algebraic minimizing controller; in particular the positive gap is algebraic and computable. This is an exact finite procedure, not a polynomial-time claim or a numerical certificate inferred from local optimization. The general procedure is proved here; a full elimination at every $M$ has not been executed in the finite diagnostics.
\end{proof}

Thus finite linear rank four does not bound the classical memory forced by a finite family of actual squared-loss tests, even on a fixed protected dimension. The number of models in this conclusion grows with the tested memory bound. It is not a claim that one fixed finite model list excludes every finite observer.

\subsection{Positive background regularizes the output face}
The leakage obstruction is sensitive to the face on which the terminal effect vanishes. Replace the effect by
$E_\eta=\eta I/2+(1-\eta)E$, $0<\eta<1$, on the same qutrit. The response becomes
\[
 u_n^{(\eta)}=\eta/2+(1-\eta)u_n.
\]
This changes a real terminal measurement, not the protected channel or the observer's information.

\begin{theorem}[Exact finite realization after background regularization]\label{thm:r70-background}
For all parameters of Theorem~\ref{thm:r70-leakage} and every $0<\eta<1$, the response $u^{(\eta)}$ has a finite exact stochastic realization. If $T\ge1$ satisfies
\begin{equation}\label{eq:r70-background-depth}
 (1-\eta)\lambda^T\le\eta/2,
\end{equation}
then at most $T+2+m(\lambda\alpha,\vartheta)$ live labels suffice. The construction uses a prescribed initial label and attains the squared-task optimum at the same cost $1/p$.

For fixed $\lambda\in(0,1)$, $\alpha=1$ and irrational phase, the least exact size tends to infinity as $\eta\downarrow0$, while the displayed upper is $O_\lambda(\log(1/\eta))$. No matching lower rate in $\eta$ is asserted. For algebraic nondegenerate inputs this finite upper and the inherited fixed-size response test give a terminating search for a minimal exact controller.
\end{theorem}
\begin{proof}
Put $a=(1-\eta)\lambda^T/2$ and $\psi=T\vartheta+\varphi$. The suffix starting at $T$ is the average of the two probability responses
\[
 f_k=\eta/2+2a\lambda^k,\qquad
 g_k=\eta/2+2a\gamma\alpha^T(\lambda\alpha)^k\cos(k\vartheta+\psi).
\]
The first has a two-label transient/absorbing realization with decoder values $\eta/2+2a$ and $\eta/2$. For the second, use the normalized eigenvector construction of Theorem~\ref{thm:r69-spectral} at eigenvalue $\lambda\alpha e^{i\vartheta}$, with decoder $\eta/2+2a\gamma\alpha^T\operatorname{Re}(e^{i\psi}v_j)$. Condition \eqref{eq:r70-background-depth} gives $2a\le\eta/2$, so all these decoder values lie in $[0,\eta]\subset[0,1]$.

Store a deterministic chain for $u_0^{(\eta)},\ldots,u_{T-1}^{(\eta)}$. At its final failure transition enter the two disjoint suffix machines with weights $1/2,1/2$. No additional initialization label is needed: the last prefix row makes and stores the choice within the counted blocks. At every later failure count the output is $(f_k+g_k)/2=u_{T+k}^{(\eta)}$. This proves the exact count and preserves the physical stopping rule.

If exact sizes stayed bounded along $\eta\downarrow0$, pad their stochastic matrices and decoders to one size and take a convergent subsequence in the compact row simplexes. Every finite coefficient is continuous in those variables. The limit would realize $u$ at $\eta=0$, contradicting Theorem~\ref{thm:r70-leakage}. The logarithmic upper follows by solving \eqref{eq:r70-background-depth} with fixed $\lambda$. For algebraic input, the target has a supplied linear representation of dimension at most four. The finite word-equality criterion of Proposition~\ref{prop:finite-realization-test}, together with this size upper, terminates a minimal exact-size search. Algebraic controller encodings and exact sampling resources remain distinct from label count.
\end{proof}
